% Part: counterfactuals
% Chapter: introduction
% Section: strict-conditional

\documentclass[../../../include/open-logic-section]{subfiles}

\begin{document}

\olfileid[hi]{cnt}{int}{str}

\olsection{कठोर सशर्त}

लुईस ने \emph{कठोर सशर्त}~$\strictif$ प्रस्तुत किया और तर्क दिया कि
निहितार्थ के अनुरूप भौतिक सशर्त नहीं, बल्कि यही है। सत्यात्मक मोडल तर्कशास्त्र
में $!A \strictif !B$ को $\Box(!A \lif !B)$ से परिभाषित किया जा सकता है।
अतः कोई कठोर सशर्त (किसी जगत पर) तभी सत्य है जब संगत भौतिक सशर्त आवश्यक हो।

भौतिक सशर्त के विरोधाभासों की तुलना में कठोर सशर्त कैसा रहता है? असत्य
पूर्ववर्ती वाला कठोर सशर्त, और सत्य परिणामांश वाला कठोर सशर्त, सत्य भी हो
सकता है और असत्य भी। इसके अतिरिक्त
$(!A \strictif !B) \lor (!B \strictif !A)$ मान्य नहीं है। कठोर सशर्त
$!A \strictif !B$, $\lnot !A \lor !B$ के समतुल्य भी नहीं है, अतः वह
सत्य-फलनीय नहीं है।

हमारे पास है:
\begin{align}
  !A \strictif !B & \Entails \lnot !A \lor !B
  \text{ किंतु:}\\
  \lnot !A \lor !B & \Entails/ !A \strictif !B\\
  !B & \Entails/ !A \strictif !B\\
  \lnot !A & \Entails/ !A \strictif !B\\
  \lnot(!A \lif !B) & \Entails/ !A \land \lnot !B
  \text{ किंतु:}\\
  !A \land \lnot !B & \Entails \lnot(!A \strictif !B)
  \intertext{फिर भी कठोर सशर्त मोडस पोनेन्स को मान्य करता है:}
  !A, !A \strictif !B & \Entails !B
  \intertext{कठोर सशर्त समुच्चित होता है:}
  !A \strictif !B,  !A \strictif !C & \Entails !A \strictif (!B \land !C)
  \intertext{कठोर सशर्त के लिए पूर्ववर्ती का प्रबलीकरण मान्य है:}
  !A \strictif !B & \Entails (!A \land !C) \strictif !B
  \intertext{कठोर सशर्त संक्रमणशील भी है:}
  !A \strictif !B, !B \strictif !C & \Entails !A \strictif !C
  \intertext{अंततः कठोर सशर्त अपने प्रतिधनात्मक के समतुल्य है:}
  !A \strictif !B & \Entails \lnot !B \strictif \lnot !A\\
  \lnot !B \strictif \lnot !A & \Entails !A \strictif !B
\end{align}

\begin{prob}
  कठोर सशर्त के लिए अमान्य निष्कर्षण-संबंधों के \Log{S5}-प्रतिदृष्टांत
  दीजिए; अर्थात निम्न के लिए:
  \begin{enumerate}
  \item $\lnot p \Entails/ \Box(p \lif q)$
  \item $q \Entails/ \Box(p \lif q)$
  \item $\lnot\Box(p \lif q) \Entails/ p \land \lnot q$
  \item $\Entails/ \Box(p \lif q) \lor \Box(q \lif p)$
  \end{enumerate}
\end{prob}

\begin{prob}
  निम्न के \Log{S5}-प्रमाण देकर दिखाइए कि मान्य निष्कर्षण-संबंध कठोर सशर्त
  के लिए सत्य हैं:
  \begin{enumerate}
  \item  $\Box(!A \lif !B) \Entails \lnot !A \lor !B$
  \item $!A \land \lnot !B \Entails \lnot\Box(!A \lif !B)$
  \item $!A, \Box(!A \lif !B) \Entails !B$
  \item $\Box(!A \lif !B), \Box(!A \lif !C) \Entails \Box(!A \lif (!B
    \land !C))$
  \item $\Box(!A \lif !B) \Entails \Box((!A \land !C) \lif !B)$
  \item $\Box(!A \lif !B), \Box(!B \lif !C) \Entails \Box(!A
    \lif !C)$
  \item $\Box(!A \lif !B) \Entails \Box(\lnot !B \lif \lnot !A)$
  \item $\Box(\lnot !B \lif \lnot !A) \Entails \Box(!A \lif !B)$
  \end{enumerate}
  \end{prob}

फिर भी कठोर सशर्त के अपने ``विरोधाभास'' हैं। जिस प्रकार असत्य पूर्ववर्ती या
सत्य परिणामांश वाला भौतिक सशर्त सत्य होता है, उसी प्रकार \emph{आवश्यकतः}
असत्य पूर्ववर्ती या आवश्यकतः सत्य परिणामांश वाला कठोर सशर्त सत्य होता है।
इसके अतिरिक्त प्रत्येक सत्य कठोर सशर्त आवश्यकतः सत्य, और प्रत्येक असत्य
कठोर सशर्त आवश्यकतः असत्य है। दूसरे शब्दों में, हमारे पास है:
\begin{align}
  \Box \lnot !A & \Entails !A \strictif !B\\
  \Box !B & \Entails !A \strictif !B\\
  !A \strictif !B& \Entails \Box(!A \strictif !B)\\
  \lnot(!A \strictif !B) & \Entails \Box\lnot(!A \strictif !B)
\end{align}
यदि $\strictif$ को ``निहित करता है'' समझें, तो ये समस्याएँ नहीं हैं। आखिर
तार्किक निष्कर्षण-संबंध गणितीय तथ्य हैं, अतः आकस्मिक नहीं हो सकते। किंतु यदि
$\strictif$ को ``यदि \dots तो \dots'' ग्रहण करने वाला तार्किक संयोजक बनाना
हो, तो ये, विशेषतः अंतिम दो, समस्याएँ उठाते हैं। निश्चय ही ``यदि \dots तो
\dots'' के कुछ कथन आकस्मिकतः सत्य या आकस्मिकतः असत्य होते हैं---वास्तव में
वे सामान्यतः न आवश्यक होते हैं, न असंभव।

\begin{prob}
  \Log{S5} में निम्न के प्रमाण दीजिए:
  \begin{enumerate}
    \item  $\Box \lnot !A \Entails !A \strictif !B$
    \item $!A \strictif !B \Entails \Box(!A \strictif !B)$
    \item $\lnot(!A \strictif !B)  \Entails \Box\lnot(!A \strictif !B)$
  \end{enumerate}
  इसके लिए $\strictif$ की परिभाषा का उपयोग कीजिए।
\end{prob}

\end{document}
